% --- Archon physics formalization source begin ---
% archon:physics
% archon:covers PhyXMiniProblems/problem_phyx_mini_0243.lean
% archon:source-report reports/phyx_mini/problem_phyx_mini_0243.source.json

\chapter{Physics problem phyx\_mini\_0243}
\label{ch:PhyXMiniProblems_problem_phyx_mini_0243}

\paragraph{Problem source.}
\#\# Physical scenario

A block of mass \textbackslash{}( M = 0.450\textbackslash{},\textbackslash{}mathrm\{kg\} \textbackslash{}) is attached to one end of a cord of mass \textbackslash{}( 0.003\textbackslash{},20\textbackslash{},\textbackslash{}mathrm\{kg\} \textbackslash{}); the other end of the cord is attached to a fixed point. The block rotates with constant angular speed in a circle on a frictionless, horizontal table as shown in figure.

\#\# Question

Through what angle does the block rotate in the time interval during which a transverse wave travels along the string from the center of the circle to the block?

\#\# Answer choices

- A: A: 0.0833rad

- B: B: 0.0863rad

- C: C: 0.0853rad

- D: D: 0.0843rad

\#\# Figure caption (auxiliary; use the image as primary evidence)

The image depicts a square table with a cube labeled "M" on top of it. The cube is attached to a string, which is secured to the table's center. The string's path creates a dashed circular line around the center of the table, indicating the possible movement path of the cube. The table is shown in a simple, isometric perspective with four legs. The cube appears to be situated at the perimeter of the dashed circle, suggesting it might move in a circular motion around the center point.

\#\# Dataset metadata

Wave Properties; Multi-Formula Reasoning, Physical Model Grounding Reasoning

\paragraph{Recorded answer/context.}
D: D: 0.0843rad

\paragraph{Figure/image path.}
/root/proposal\_for\_physic/science-mango/phyx\_mini\_run/phyx\_data/test\_image/243.png

\paragraph{Formalization target.}
create a compiling Lean file with sorry bodies at `PhyXMiniProblems/problem\_phyx\_mini\_0243.lean`. The Lean declarations must preserve the physical quantities, dimensions or dimensional roles, figure labels, governing-law hypotheses, and final relation expressed by this problem.
Use Mathlib/Physlib names found through LeanExplore where available. If a physics API is missing, introduce faithful local abstractions rather than scalar placeholder aliases.

\begin{theorem}[Physics formalization target]
\label{thm:physics:phyx_mini_0243:target}
\lean{PhyXMiniProblems.ProblemPhyXMini0243.angleDuringWaveTravel_matches_recordedAnswerD}
\uses{def:physics:phyx-mini-0243:phyxminiproblems-problemphyxmini0243-rotatingblockcordexperiment, def:physics:phyx-mini-0243:phyxminiproblems-problemphyxmini0243-angleduringwavetravelradians, def:physics:phyx-mini-0243:phyxminiproblems-problemphyxmini0243-matchesproblemandprimaryfigure, def:physics:phyx-mini-0243:phyxminiproblems-problemphyxmini0243-hasphysicalparameters, def:physics:phyx-mini-0243:phyxminiproblems-problemphyxmini0243-satisfiesrotatingcordwavelaws, def:physics:phyx-mini-0243:phyxminiproblems-problemphyxmini0243-recordedanswerchoice, def:physics:phyx-mini-0243:phyxminiproblems-problemphyxmini0243-matchesdisplayedanswer}
The assigned autoformalize agent should translate this physics problem into Lean declarations in the covered file, with theorem and lemma proof bodies written as `by sorry`.
\end{theorem}
\begin{proof}
This is an autoformalization task, not a proof task. Produce statements that can later be proved by the physics prover without weakening the physical model.
\end{proof}

% --- Archon declaration topology begin ---
\section{Declaration topology}
\begin{definition}[Mass Quantity]
  \label{def:physics:phyx-mini-0243:phyxminiproblems-problemphyxmini0243-massquantity}
  \lean{PhyXMiniProblems.ProblemPhyXMini0243.MassQuantity}
  A nonnegative physical mass
\end{definition}
\begin{proof}
  This is the declared model definition of Mass Quantity.
\end{proof}
\begin{definition}[Length Quantity]
  \label{def:physics:phyx-mini-0243:phyxminiproblems-problemphyxmini0243-lengthquantity}
  \lean{PhyXMiniProblems.ProblemPhyXMini0243.LengthQuantity}
  A nonnegative physical length
\end{definition}
\begin{proof}
  This is the declared model definition of Length Quantity.
\end{proof}
\begin{definition}[Time Quantity]
  \label{def:physics:phyx-mini-0243:phyxminiproblems-problemphyxmini0243-timequantity}
  \lean{PhyXMiniProblems.ProblemPhyXMini0243.TimeQuantity}
  A nonnegative physical duration
\end{definition}
\begin{proof}
  This is the declared model definition of Time Quantity.
\end{proof}
\begin{definition}[Angular Speed Quantity]
  \label{def:physics:phyx-mini-0243:phyxminiproblems-problemphyxmini0243-angularspeedquantity}
  \lean{PhyXMiniProblems.ProblemPhyXMini0243.AngularSpeedQuantity}
  A nonnegative angular speed; radians are dimensionless
\end{definition}
\begin{proof}
  This is the declared model definition of Angular Speed Quantity.
\end{proof}
\begin{definition}[Tension Quantity]
  \label{def:physics:phyx-mini-0243:phyxminiproblems-problemphyxmini0243-tensionquantity}
  \lean{PhyXMiniProblems.ProblemPhyXMini0243.TensionQuantity}
  A nonnegative force, used for the tension in the cord
\end{definition}
\begin{proof}
  This is the declared model definition of Tension Quantity.
\end{proof}
\begin{definition}[Linear Mass Density Quantity]
  \label{def:physics:phyx-mini-0243:phyxminiproblems-problemphyxmini0243-linearmassdensityquantity}
  \lean{PhyXMiniProblems.ProblemPhyXMini0243.LinearMassDensityQuantity}
  A nonnegative mass per unit length of cord
\end{definition}
\begin{proof}
  This is the declared model definition of Linear Mass Density Quantity.
\end{proof}
\begin{definition}[Speed Quantity]
  \label{def:physics:phyx-mini-0243:phyxminiproblems-problemphyxmini0243-speedquantity}
  \lean{PhyXMiniProblems.ProblemPhyXMini0243.SpeedQuantity}
  A nonnegative propagation speed along the cord
\end{definition}
\begin{proof}
  This is the declared model definition of Speed Quantity.
\end{proof}
\begin{definition}[mass In Kilograms]
  \label{def:physics:phyx-mini-0243:phyxminiproblems-problemphyxmini0243-massinkilograms}
  \lean{PhyXMiniProblems.ProblemPhyXMini0243.massInKilograms}
  \uses{def:physics:phyx-mini-0243:phyxminiproblems-problemphyxmini0243-massquantity}
  Kilogram readout of a mass in coherent SI units
\end{definition}
\begin{proof}
  This is the declared model definition of mass In Kilograms.
\end{proof}
\begin{definition}[length In Meters]
  \label{def:physics:phyx-mini-0243:phyxminiproblems-problemphyxmini0243-lengthinmeters}
  \lean{PhyXMiniProblems.ProblemPhyXMini0243.lengthInMeters}
  \uses{def:physics:phyx-mini-0243:phyxminiproblems-problemphyxmini0243-lengthquantity}
  Meter readout of a length in coherent SI units
\end{definition}
\begin{proof}
  This is the declared model definition of length In Meters.
\end{proof}
\begin{definition}[time In Seconds]
  \label{def:physics:phyx-mini-0243:phyxminiproblems-problemphyxmini0243-timeinseconds}
  \lean{PhyXMiniProblems.ProblemPhyXMini0243.timeInSeconds}
  \uses{def:physics:phyx-mini-0243:phyxminiproblems-problemphyxmini0243-timequantity}
  Second readout of a duration in coherent SI units
\end{definition}
\begin{proof}
  This is the declared model definition of time In Seconds.
\end{proof}
\begin{definition}[angular Speed In Radians Per Second]
  \label{def:physics:phyx-mini-0243:phyxminiproblems-problemphyxmini0243-angularspeedinradianspersecond}
  \lean{PhyXMiniProblems.ProblemPhyXMini0243.angularSpeedInRadiansPerSecond}
  \uses{def:physics:phyx-mini-0243:phyxminiproblems-problemphyxmini0243-angularspeedquantity}
  Radian-per-second readout of an angular speed
\end{definition}
\begin{proof}
  This is the declared model definition of angular Speed In Radians Per Second.
\end{proof}
\begin{definition}[tension In Newtons]
  \label{def:physics:phyx-mini-0243:phyxminiproblems-problemphyxmini0243-tensioninnewtons}
  \lean{PhyXMiniProblems.ProblemPhyXMini0243.tensionInNewtons}
  \uses{def:physics:phyx-mini-0243:phyxminiproblems-problemphyxmini0243-tensionquantity}
  Newton readout of a cord tension
\end{definition}
\begin{proof}
  This is the declared model definition of tension In Newtons.
\end{proof}
\begin{definition}[linear Mass Density In Kilograms Per Meter]
  \label{def:physics:phyx-mini-0243:phyxminiproblems-problemphyxmini0243-linearmassdensityinkilogramspermeter}
  \lean{PhyXMiniProblems.ProblemPhyXMini0243.linearMassDensityInKilogramsPerMeter}
  \uses{def:physics:phyx-mini-0243:phyxminiproblems-problemphyxmini0243-linearmassdensityquantity}
  Kilogram-per-meter readout of a linear mass density
\end{definition}
\begin{proof}
  This is the declared model definition of linear Mass Density In Kilograms Per Meter.
\end{proof}
\begin{definition}[speed In Meters Per Second]
  \label{def:physics:phyx-mini-0243:phyxminiproblems-problemphyxmini0243-speedinmeterspersecond}
  \lean{PhyXMiniProblems.ProblemPhyXMini0243.speedInMetersPerSecond}
  \uses{def:physics:phyx-mini-0243:phyxminiproblems-problemphyxmini0243-speedquantity}
  Meter-per-second readout of a wave speed
\end{definition}
\begin{proof}
  This is the declared model definition of speed In Meters Per Second.
\end{proof}
\begin{definition}[Cord End]
  \label{def:physics:phyx-mini-0243:phyxminiproblems-problemphyxmini0243-cordend}
  \lean{PhyXMiniProblems.ProblemPhyXMini0243.CordEnd}
  The two physically distinct ends of the cord
\end{definition}
\begin{proof}
  This is the declared model definition of Cord End.
\end{proof}
\begin{definition}[Figure Object]
  \label{def:physics:phyx-mini-0243:phyxminiproblems-problemphyxmini0243-figureobject}
  \lean{PhyXMiniProblems.ProblemPhyXMini0243.FigureObject}
  Object labels visible or implied in the primary figure
\end{definition}
\begin{proof}
  This is the declared model definition of Figure Object.
\end{proof}
\begin{definition}[Orbit Shape]
  \label{def:physics:phyx-mini-0243:phyxminiproblems-problemphyxmini0243-orbitshape}
  \lean{PhyXMiniProblems.ProblemPhyXMini0243.OrbitShape}
  The dashed trajectory drawn on the tabletop
\end{definition}
\begin{proof}
  This is the declared model definition of Orbit Shape.
\end{proof}
\begin{definition}[Table Orientation]
  \label{def:physics:phyx-mini-0243:phyxminiproblems-problemphyxmini0243-tableorientation}
  \lean{PhyXMiniProblems.ProblemPhyXMini0243.TableOrientation}
  Orientation of the surface supporting the block
\end{definition}
\begin{proof}
  This is the declared model definition of Table Orientation.
\end{proof}
\begin{definition}[Surface Condition]
  \label{def:physics:phyx-mini-0243:phyxminiproblems-problemphyxmini0243-surfacecondition}
  \lean{PhyXMiniProblems.ProblemPhyXMini0243.SurfaceCondition}
  Contact model for the table surface
\end{definition}
\begin{proof}
  This is the declared model definition of Surface Condition.
\end{proof}
\begin{definition}[Rotating Block Cord Experiment]
  \label{def:physics:phyx-mini-0243:phyxminiproblems-problemphyxmini0243-rotatingblockcordexperiment}
  \lean{PhyXMiniProblems.ProblemPhyXMini0243.RotatingBlockCordExperiment}
  \uses{def:physics:phyx-mini-0243:phyxminiproblems-problemphyxmini0243-massquantity, def:physics:phyx-mini-0243:phyxminiproblems-problemphyxmini0243-lengthquantity, def:physics:phyx-mini-0243:phyxminiproblems-problemphyxmini0243-timequantity, def:physics:phyx-mini-0243:phyxminiproblems-problemphyxmini0243-angularspeedquantity, def:physics:phyx-mini-0243:phyxminiproblems-problemphyxmini0243-tensionquantity, def:physics:phyx-mini-0243:phyxminiproblems-problemphyxmini0243-linearmassdensityquantity, def:physics:phyx-mini-0243:phyxminiproblems-problemphyxmini0243-speedquantity, def:physics:phyx-mini-0243:phyxminiproblems-problemphyxmini0243-cordend, def:physics:phyx-mini-0243:phyxminiproblems-problemphyxmini0243-figureobject, def:physics:phyx-mini-0243:phyxminiproblems-problemphyxmini0243-orbitshape, def:physics:phyx-mini-0243:phyxminiproblems-problemphyxmini0243-tableorientation, def:physics:phyx-mini-0243:phyxminiproblems-problemphyxmini0243-surfacecondition}
  The independent physical quantities and figure-derived roles in the setup. angularDisplacementRadians duration is an unwrapped scalar radian readout, appropriate for the short time intervals in this problem. It is not defined to be the requested answer; its relation to the angular speed is stated later as a governing kinematic law
\end{definition}
\begin{proof}
  This is the declared model definition of Rotating Block Cord Experiment.
\end{proof}
\begin{definition}[angle During Wave Travel Radians]
  \label{def:physics:phyx-mini-0243:phyxminiproblems-problemphyxmini0243-angleduringwavetravelradians}
  \lean{PhyXMiniProblems.ProblemPhyXMini0243.angleDuringWaveTravelRadians}
  \uses{def:physics:phyx-mini-0243:phyxminiproblems-problemphyxmini0243-rotatingblockcordexperiment}
  The angle swept by the block during the center-to-block wave transit
\end{definition}
\begin{proof}
  This is the declared model definition of angle During Wave Travel Radians.
\end{proof}
\begin{definition}[Matches Problem And Primary Figure]
  \label{def:physics:phyx-mini-0243:phyxminiproblems-problemphyxmini0243-matchesproblemandprimaryfigure}
  \lean{PhyXMiniProblems.ProblemPhyXMini0243.MatchesProblemAndPrimaryFigure}
  \uses{def:physics:phyx-mini-0243:phyxminiproblems-problemphyxmini0243-massinkilograms, def:physics:phyx-mini-0243:phyxminiproblems-problemphyxmini0243-rotatingblockcordexperiment}
  Numerical data and geometry stated by the problem or read from the primary figure. The source's 0.003,20 kg typography is normalized to 0.00320 kg
\end{definition}
\begin{proof}
  This is the declared model definition of Matches Problem And Primary Figure.
\end{proof}
\begin{definition}[Has Physical Parameters]
  \label{def:physics:phyx-mini-0243:phyxminiproblems-problemphyxmini0243-hasphysicalparameters}
  \lean{PhyXMiniProblems.ProblemPhyXMini0243.HasPhysicalParameters}
  \uses{def:physics:phyx-mini-0243:phyxminiproblems-problemphyxmini0243-massinkilograms, def:physics:phyx-mini-0243:phyxminiproblems-problemphyxmini0243-lengthinmeters, def:physics:phyx-mini-0243:phyxminiproblems-problemphyxmini0243-timeinseconds, def:physics:phyx-mini-0243:phyxminiproblems-problemphyxmini0243-angularspeedinradianspersecond, def:physics:phyx-mini-0243:phyxminiproblems-problemphyxmini0243-tensioninnewtons, def:physics:phyx-mini-0243:phyxminiproblems-problemphyxmini0243-linearmassdensityinkilogramspermeter, def:physics:phyx-mini-0243:phyxminiproblems-problemphyxmini0243-speedinmeterspersecond, def:physics:phyx-mini-0243:phyxminiproblems-problemphyxmini0243-rotatingblockcordexperiment}
  Positivity and nondegeneracy conditions for the physical setup
\end{definition}
\begin{proof}
  This is the declared model definition of Has Physical Parameters.
\end{proof}
\begin{definition}[Satisfies Rotating Cord Wave Laws]
  \label{def:physics:phyx-mini-0243:phyxminiproblems-problemphyxmini0243-satisfiesrotatingcordwavelaws}
  \lean{PhyXMiniProblems.ProblemPhyXMini0243.SatisfiesRotatingCordWaveLaws}
  \uses{def:physics:phyx-mini-0243:phyxminiproblems-problemphyxmini0243-timequantity, def:physics:phyx-mini-0243:phyxminiproblems-problemphyxmini0243-massinkilograms, def:physics:phyx-mini-0243:phyxminiproblems-problemphyxmini0243-lengthinmeters, def:physics:phyx-mini-0243:phyxminiproblems-problemphyxmini0243-timeinseconds, def:physics:phyx-mini-0243:phyxminiproblems-problemphyxmini0243-angularspeedinradianspersecond, def:physics:phyx-mini-0243:phyxminiproblems-problemphyxmini0243-tensioninnewtons, def:physics:phyx-mini-0243:phyxminiproblems-problemphyxmini0243-linearmassdensityinkilogramspermeter, def:physics:phyx-mini-0243:phyxminiproblems-problemphyxmini0243-speedinmeterspersecond, def:physics:phyx-mini-0243:phyxminiproblems-problemphyxmini0243-rotatingblockcordexperiment}
  The governing textbook laws used for this question: uniform linear density mu = m / R; the light-cord, uniform-tension centripetal relation T = M * omega\textasciicircum{}2 * R; transverse wave speed v = sqrt (T / mu); center-to-block transit time t = R / v; constant-angular-speed kinematics theta = omega * t. These laws relate independently stored physical quantities. In particular, they do not contain either sqrt (m / M) or the displayed answer 0.0843
\end{definition}
\begin{proof}
  This is the declared model definition of Satisfies Rotating Cord Wave Laws.
\end{proof}
\begin{lemma}[angle During Wave Travel mass Ratio]
  \label{lem:physics:phyx-mini-0243:phyxminiproblems-problemphyxmini0243-angleduringwavetravel-massratio}
  \lean{PhyXMiniProblems.ProblemPhyXMini0243.angleDuringWaveTravel_massRatio}
  \uses{def:physics:phyx-mini-0243:phyxminiproblems-problemphyxmini0243-massinkilograms, def:physics:phyx-mini-0243:phyxminiproblems-problemphyxmini0243-rotatingblockcordexperiment, def:physics:phyx-mini-0243:phyxminiproblems-problemphyxmini0243-angleduringwavetravelradians, def:physics:phyx-mini-0243:phyxminiproblems-problemphyxmini0243-matchesproblemandprimaryfigure, def:physics:phyx-mini-0243:phyxminiproblems-problemphyxmini0243-hasphysicalparameters, def:physics:phyx-mini-0243:phyxminiproblems-problemphyxmini0243-satisfiesrotatingcordwavelaws}
  The governing laws imply that the cord length and angular speed cancel, so the swept angle is the square root of the cord-to-block mass ratio. This is a derived conclusion, not a law or data field
\end{lemma}
\begin{proof}
  The conclusion follows from the governing relations and figure readouts named in the statement of angle During Wave Travel mass Ratio.
\end{proof}
\begin{definition}[Answer Choice]
  \label{def:physics:phyx-mini-0243:phyxminiproblems-problemphyxmini0243-answerchoice}
  \lean{PhyXMiniProblems.ProblemPhyXMini0243.AnswerChoice}
  Labels printed beside the four candidate angles
\end{definition}
\begin{proof}
  This is the declared model definition of Answer Choice.
\end{proof}
\begin{definition}[radians]
  \label{def:physics:phyx-mini-0243:phyxminiproblems-problemphyxmini0243-answerchoice-radians}
  \lean{PhyXMiniProblems.ProblemPhyXMini0243.AnswerChoice.radians}
  \uses{def:physics:phyx-mini-0243:phyxminiproblems-problemphyxmini0243-answerchoice}
  Unwrapped radian value printed beside each answer label
\end{definition}
\begin{proof}
  This is the declared model definition of radians.
\end{proof}
\begin{definition}[recorded Answer Choice]
  \label{def:physics:phyx-mini-0243:phyxminiproblems-problemphyxmini0243-recordedanswerchoice}
  \lean{PhyXMiniProblems.ProblemPhyXMini0243.recordedAnswerChoice}
  \uses{def:physics:phyx-mini-0243:phyxminiproblems-problemphyxmini0243-answerchoice}
  The answer label recorded in the supplied dataset
\end{definition}
\begin{proof}
  This is the declared model definition of recorded Answer Choice.
\end{proof}
\begin{definition}[Matches Displayed Answer]
  \label{def:physics:phyx-mini-0243:phyxminiproblems-problemphyxmini0243-matchesdisplayedanswer}
  \lean{PhyXMiniProblems.ProblemPhyXMini0243.MatchesDisplayedAnswer}
  \uses{def:physics:phyx-mini-0243:phyxminiproblems-problemphyxmini0243-answerchoice}
  Agreement with a four-decimal-place displayed angle: the exact radian value is within half a unit in the last displayed decimal place
\end{definition}
\begin{proof}
  This is the declared model definition of Matches Displayed Answer.
\end{proof}
% --- Archon declaration topology end ---
% --- Archon physics formalization source end ---
